\subsectionoptional{Jordan 标准形}
我们要寻找的是矩阵相似之下的一个标准形。本小节通过从幂零矩阵各相似类的标准形推进到全体相似类的标准形，来完成这一工作。

\begin{lemma} \label{le:NilIffOnlyEigenZero}
非平凡向量空间上的线性变换是幂零的，当且仅当它仅有的特征值为零。
\end{lemma}

\begin{proof}
若非平凡向量空间上的线性变换 $t$ 是幂零的，则存在 $n$ 使得 $t^n$ 是零映射，从而 $t$ 满足多项式 $p(x)=x^n=(x-0)^n$。由 \nearbylemma{le:tSatisImpMinPolyDivides}，$t$ 的极小多项式整除 $p$，故该极小多项式仅有的根是零。再由 Cayley-Hamilton 定理，\nearbytheorem{th:CayHam}，特征多项式仅有的根是零。于是 $t$ 仅有的特征值为零。

反之，若 \( n \) 维空间上的变换 \( t \) 仅有的特征值是零，则它的特征多项式是 \( x^n \)。\nearbylemma{le:MatSatItsCharPoly} 指出映射满足其特征多项式，故 \( t^n \) 是零映射。于是 $t$ 是幂零的。
\end{proof}

\noindent 该引理的陈述中出现“非平凡向量空间”，是因为在平凡空间 $\set{\zero}$ 上唯一的变换是零映射，而它没有相关的非零特征向量，因此没有特征值。

\begin{corollary} \label{cor:tMinLambdaNilpotent}
变换 $t-\lambda$ 是幂零的，当且仅当 $t$ 仅有的特征值是 $\lambda$。
\end{corollary}

\begin{proof}
变换 \( t-\lambda \) 是幂零的，当且仅当 $t-\lambda$ 仅有的特征值是 \( 0 \)。这当且仅当 $t$ 仅有的特征值是 $\lambda$ 时才成立，因为 \( t(\vec{v})=\lambda\vec{v} \,\) 当且仅当 \( (t-\lambda)\,(\vec{v})=0\cdot\vec{v} \)。
\end{proof}

\begin{lemma}   \label{le:SimRespAddScalar}
若矩阵 \( T-\lambda I \) 与 \( N \) 相似，则 \( T \) 与 \( N+\lambda I \) 也相似，且经由相同的换基矩阵。
\end{lemma}

\begin{proof}
由 \( N=P(T-\lambda I)P^{-1}=PTP^{-1}-P(\lambda I)P^{-1} \)，我们得到 $N=PTP^{-1}-PP^{-1}(\lambda I)$，因为对角矩阵 \( \lambda I \) 与任何矩阵交换。于是 \( N=PTP^{-1}-\lambda I \)，从而 \( N+\lambda I=PTP^{-1} \)。
\end{proof}
对于幂零矩阵的情形，也就是仅有一个特征值为零的矩阵的情形，我们已经有其标准形。由推论~III.\ref{cor:NilpotentMatCanonForm}，每个这样的矩阵都相似于一个除了由次对角线上的~1 构成的块以外全为零的矩阵。
% (To make this representation unique we can fix some arrangement of
% the blocks, perhaps from longest to shortest.)
前面的结论使我们能把它推广到仅有一个特征值~$\lambda$ 的所有矩阵。在新形式中，除了由次对角线上的~1 构成的块以外，对角线上的所有元素都是~$\lambda$。
（见下文的 \nearbydefinition{def:JordanBlock}。）

\begin{example}   \label{ex:SingJordBlock}
矩阵
\begin{equation*}
  T=\begin{mat}[r]
      2  &-1  \\
      1  &4
    \end{mat}
\end{equation*}
的特征多项式是 \( (x-3)^2 \)，因此 \( T \) 仅有一个特征值 \( 3 \)。
% We can find a matrix similar to $T$ that is simple in
% the way that our canonical form for nilpotent matrices is simple.
% The form is described in \nearbydefinition{def:JordanBlock} below but
% as a preview here the new matrix.
% \begin{equation*}
%   \begin{mat}[r]
%       3  &0  \\
%       1  &3
%     \end{mat}
% \end{equation*}

由于 $3$ 是仅有的特征值，$T-3I$ 仅有的特征值是 \( 0 \)，从而是幂零的。为求~$T$ 的标准形，先按前一小节的做法为 $T-3I$ 求出一个链基：计算该矩阵的各次幂，并求出这些幂的零空间。
\begin{center}
  \begin{tabular}{c|cc}
    \multicolumn{1}{c}{\textit{幂次} \( p \)}  &\( (T-3I)^p \)  &\( \nullspace{(T-3I)^p}  \)   \\
    \hline
    \( 1 \)
    &\(
       \matrixvenlarge{\begin{mat}[r]
         -1  &-1 \\
         1   &1   \\
       \end{mat}}  \)
    &\( \set{\matrixvenlarge{\colvec{-y \\ y}}\suchthat
                               y\in\C}  \)   \\
    \( 2 \)
    &\(
       \matrixvenlarge{\begin{mat}[r]
         0  &0  \\
         0  &0   \\
       \end{mat}}  \)
    &\( \set{\matrixvenlarge{\colvec{x \\ y}}\suchthat
                               x,y\in\C}  \)   \\
  \end{tabular}
\end{center}
$t-3$ 的零空间维数为一，$(t-3)^2$ 的零空间维数为二。因此，下面就是标准表示矩阵，以及 $t-3$ 相应的链基~$B$ 的形式。
\begin{equation*}
  \rep{t-3}{B,B}
  =N=
  \begin{mat}[r]
      0  &0   \\
      1  &0
  \end{mat}
  \qquad
  \begin{array}{r}
    \vec{\beta}_1\xmapsto{t-3}\vec{\beta}_2\xmapsto{t-3}\zero  \\
  \end{array}
\end{equation*}
求这样一个基的方法是：先取一个被 $t-3$ 映到~$\zero$ 的 $\vec{\beta}_2$（也就是从表中给出的零空间 $\nullspace{t-3}$ 中取一个元素）。再选取一个被 $(t-3)^2$ 映到~$\zero$ 的 $\vec{\beta}_1$。这一选取除了要求属于零空间以外，唯一的限制是 $\vec{\beta}_1$ 与 $\vec{\beta}_2$ 必须构成线性无关集。下面是一种可能的取法。
\begin{equation*}
  \vec{\beta}_1=\colvec[r]{1 \\ 1}
  \quad
  \vec{\beta}_2=\colvec[r]{-2 \\ 2}
\end{equation*}
由此，\nearbylemma{le:SimRespAddScalar} 表明 \( T \) 相似于这个矩阵。
\begin{equation*}
  \rep{t}{B,B}=
  N+3I=
  \begin{mat}[r]
     3  &0  \\
     1  &3
  \end{mat}
\end{equation*}

我们可以把相似计算明确地展示出来。取 $T$ 为变换 $\map{t}{\C^2}{\C^2}$ 相对于标准基的表示矩阵（本章余下部分都将如此）。
% Recall how to find the change of
% basis matrices $P$ and $P^{-1}$ to express \( N \) as \( P(T-3I)P^{-1} \).
相似示意图
\begin{equation*}
  \begin{CD}
    \C^2_{\wrt{\stdbasis_2}}      @>t-3>T-3I>      \C^2_{\wrt{\stdbasis_2}}     \\
    @V\scriptstyle\identity V\scriptstyle PV
                                 @V\scriptstyle\identity V\scriptstyle PV \\
    \C^2_{\wrt{B}}                 @>t-3>N>         \C^2_{\wrt{B}}
  \end{CD}
\end{equation*}
说明从左下角走到左上角要乘以
\begin{equation*}
  P^{-1}=\bigl(\rep{\identity}{\stdbasis_2,B}\bigr)^{-1}
    =\rep{\identity}{B,\stdbasis_2}
    =\begin{mat}[r]
        1  &-2  \\
        1  &2
     \end{mat}
\end{equation*}
而从右上角走到右下角则要乘以这个矩阵。
\begin{equation*}
  P=\begin{mat}[r]
      1  &-2  \\
      1  &2
     \end{mat}^{-1}\!\!
   =\begin{mat}[r]
      1/2  &1/2  \\
      -1/4 &1/4
   \end{mat}
\end{equation*}
于是这个等式展示了相似性。
\begin{equation*}
  \begin{mat}[r]
      1/2  &1/2  \\
      -1/4 &1/4
   \end{mat}
   \begin{mat}[r]
      2  &-1  \\
      1  &4
    \end{mat}
    \begin{mat}[r]
        1  &-2  \\
        1  &2
     \end{mat}
  =
  \begin{mat}[r]
     3  &0  \\
     1  &3
  \end{mat}
\end{equation*}
\end{example}

\begin{example}  \label{ex:FullNilpotentCase}
矩阵
\begin{equation*}
  T=
  \begin{mat}[r]
    4  &1  &0  &-1  \\
    0  &3  &0  &1   \\
    0  &0  &4  &0   \\
    1  &0  &0  &5
   \end{mat}
\end{equation*}
的特征多项式是 \( (x-4)^4 \)，
因此仅有一个特征值~$4$。
\begin{center}
  \begin{tabular}{c|cc}
    \multicolumn{1}{c}{\textit{幂次} \( p \)}  &\( (T-4I)^p \)  &\( \nullspace{(T-4I)^p}  \)   \\
    \hline
    \( 1 \)
    &\(
       \matrixvenlarge{\begin{mat}[r]
         0  &1  &0  &-1  \\
         0  &-1 &0  &1   \\
         0  &0  &0  &0   \\
         1  &0  &0  &1
       \end{mat}}  \)
    &\( \set{\matrixvenlarge{\colvec{-w \\ w \\ z \\ w}}\suchthat
                               z,w\in\C}  \)   \\
    \( 2 \)
    &\(
       \matrixvenlarge{\begin{mat}[r]
         -1  &-1 &0  &0  \\
          1  &1  &0  &0   \\
          0  &0  &0  &0   \\
          1  &1  &0  &0
       \end{mat}}  \)
    &\( \set{\matrixvenlarge{\colvec{-y \\ y \\ z \\ w}}\suchthat
                               y,z,w\in\C}  \)   \\
    \( 3 \)
    &\(
       \matrixvenlarge{\begin{mat}[r]
          0  &0   &0  &0  \\
          0  &0  &0  &0   \\
          0  &0  &0  &0   \\
          0  &0  &0  &0   \\
       \end{mat}}  \)
    &\( \set{\matrixvenlarge{\colvec{x \\ y \\ z \\ w}}\suchthat
                               x,y,z,w\in\C}  \)   \\
  \end{tabular}
% sage: T=matrix(QQ, [[4,1,0,-1], [0,3,0,1], [0,0,4,0], [1,0,0,5]])
% sage: S=T-4*identity_matrix(4)
% sage: S^2
% [-1 -1  0  0]
% [ 1  1  0  0]
% [ 0  0  0  0]
% [ 1  1  0  0]
% sage: S^3
% [0 0 0 0]
% [0 0 0 0]
% [0 0 0 0]
% [0 0 0 0]
% sage: S
% [ 0  1  0 -1]
% [ 0 -1  0  1]
% [ 0  0  0  0]
% [ 1  0  0  1]
% sage: S*vector(QQ, [-1,1, 0, 1])
% (0, 0, 0, 0)
\end{center}
$t-4$ 的零空间维数为二，$(t-4)^2$
的零空间维数为三，而 $(t-4)^3$ 的零空间维数为四。
由此得到 $t-4$ 的标准形。
\begin{equation*}
  N=\rep{t-4}{B,B}=
  \begin{mat}[r]
    0  &0  &0  &0   \\
    1  &0  &0  &0   \\
    0  &1  &0  &0   \\
    0  &0  &0  &0
   \end{mat}
   \qquad
  \begin{array}{r}
    \vec{\beta}_{1}\xmapsto{t-4}\vec{\beta}_{2}
       \xmapsto{t-4}\vec{\beta}_{3}\xmapsto{t-4}\zero \\
   \vec{\beta}_{4}\xmapsto{t-4}\zero \\
  \end{array}
\end{equation*}
为构造这样一个基，选取被~$t-4$ 映到~$\zero$
的 $\vec{\beta}_3$ 和~$\vec{\beta}_4$。
再选取一个被~$(t-4)^2$ 映到~$\zero$ 的~$\vec{\beta}_2$，
以及一个被~$(t-4)^3$ 映到~$\zero$ 的 $\vec{\beta}_1$。
\begin{equation*}
  \vec{\beta}_1=\colvec{1 \\ 0 \\ 0 \\ 0}
  \quad
  \vec{\beta}_2=\colvec{-1 \\ 1 \\ 0 \\ 0}
  \quad
  \vec{\beta}_3=\colvec{-1 \\ 1 \\ 0 \\ 1}
  \quad
  \vec{\beta}_4=\colvec{0 \\ 0 \\ 1 \\ 0}
\end{equation*}
注意，选取这四个向量时避免了任何线性相关。
这就是与~$T$ 相似的标准形矩阵。
\begin{equation*}
  \rep{t}{B,B}
  =N+4I=
  \begin{mat}[r]
    4  &0  &0  &0   \\
    1  &4  &0  &0   \\
    0  &1  &4  &0   \\
    0  &0  &0  &4
   \end{mat}
\end{equation*}
\end{example}

\begin{definition}  \label{def:JordanBlock}
\definend{Jordan 块}\index{Jordan block}
是一个方阵，或者矩阵内部的一个方形元素块，
它除以下情形外全为零：存在一个数 $\lambda\in\C$，
使得每个对角元都是 $\lambda$，
且每个次对角线上的元素都是~$1$。
（若该块是 $\nbyn{1}$ 的，则它没有次对角线上的~1。）
\end{definition}

相应基中的这些链称为
\definend{Jordan 链}\index{Jordan string}\index{basis!Jordan string}
或 \definend{Jordan 链}。\index{Jordan chain}\index{basis!Jordan chain}

上面的例子说明，对于仅有一个特征值的矩阵，
Jordan 块矩阵是各相似类的标准形代表元。
我们可以把这种矩阵形式变得唯一，做法是安排基元素的次序，
使得由次对角线上的 1 构成的块按从左到右由长到短排列。

\begin{example}
仅有的特征值为 \( 1/2 \) 的 \( \nbyn{3} \) 矩阵分成
三个相似类。
这三个类具有如下标准形代表元。
\begin{equation*}
  \begin{mat}[r]
     1/2  &0    &0  \\
     0    &1/2  &0  \\
     0    &0    &1/2
   \end{mat}
   \qquad
   \begin{mat}[r]
     1/2  &0    &0  \\
     1    &1/2  &0  \\
     0    &0    &1/2
   \end{mat}
   \qquad
   \begin{mat}[r]
     1/2  &0    &0  \\
     1    &1/2  &0  \\
     0    &1    &1/2
   \end{mat}
\end{equation*}
特别地，矩阵
\begin{equation*}
   \begin{mat}[r]
     1/2  &0    &0    \\
     0    &1/2  &0    \\
     0    &1    &1/2
   \end{mat}
\end{equation*}
属于由中间那个矩阵代表的相似类，这是由
次对角线上的 1 构成的块的排列约定所致。
\end{example}

现在我们已准备好完成本章的计划。
首先回顾一下到目前为止我们关于
变换 $\map{t}{\C^n}{\C^n}$ 的所得。

可对角化的情形是指
$t$ 有~$n$ 个互不相同的特征值
$\lambda_1$, \ldots, $\lambda_n$，也就是特征值的个数
等于空间的维数的情形。
此时有一组基
$\sequence{\vec{\beta}_1,\ldots,\vec{\beta}_n}$，其中
$\vec{\beta}_i$ 是与特征值~$\lambda_i$ 相伴的特征向量。
下例中
$n=3$，三个特征值为 $\lambda_1=1$、$\lambda_2=3$
和 $\lambda_3=-1$。
图中给出了标准形代表矩阵及相应的基。
\begin{equation*}
  T_1=\begin{mat}
    1  &0  &0  \\
    0  &3  &0  \\
    0  &0  &-1
  \end{mat}
  \qquad
    \begin{array}{r}
    \vec{\beta}_1\mapsunder{t-1}\zero  \\
    \vec{\beta}_2\mapsunder{t-3}\zero  \\
    \vec{\beta}_3\mapsunder{t+1}\zero
    \end{array}
\end{equation*}
对角化的一个例子是
Five.II.\ref{ex:SummaryOfDiagonalizable}。

$t$ 仅有一个特征值的情形则利用幂零性的结果，
得到由构成不相交链的
相伴特征向量组成的基。
本例有 $n=10$，唯一的特征值~$\lambda=2$，
以及五个 Jordan 块。
\begin{equation*}
  T_2=\begin{mat}
    2  &0  &0  &0 &0 &0 &0 &0 &0 &0 \\
    1  &2  &0  &0 &0 &0 &0 &0 &0 &0 \\
    0  &1  &2  &0 &0 &0 &0 &0 &0 &0 \\
    0  &0  &0  &2 &0 &0 &0 &0 &0 &0 \\
    0  &0  &0  &1 &2 &0 &0 &0 &0 &0 \\
    0  &0  &0  &0 &1 &2 &0 &0 &0 &0 \\
    0  &0  &0  &0 &0 &0 &2 &0 &0 &0 \\
    0  &0  &0  &0 &0 &0 &1 &2 &0 &0 \\
    0  &0  &0  &0 &0 &0 &0 &0 &2 &0 \\
    0  &0  &0  &0 &0 &0 &0 &0 &0 &2 \\
  \end{mat}
  \qquad
  {\renewcommand{\arraystretch}{1.75}
    \begin{array}{r}
    \vec{\beta}_1\mapsunder{t-2}\vec{\beta}_2\mapsunder{t-2}\vec{\beta}_3\mapsunder{t-2}\zero  \\
    \vec{\beta}_4\mapsunder{t-2}\vec{\beta}_5\mapsunder{t-2}\vec{\beta}_6\mapsunder{t-2}\zero  \\
    \vec{\beta}_7\mapsunder{t-2}\vec{\beta}_8\mapsunder{t-2}\zero  \\
    \vec{\beta}_9\mapsunder{t-2}\zero  \\
    \vec{\beta}_{10}\mapsunder{t-2}\zero
    \end{array}}
\end{equation*}
我们已在 \nearbyexample{ex:FullNilpotentCase}
看到一个完整的例子。

下面的 \nearbytheorem{th:JordanForm} 把这两种情形
推广到任意变换
$\map{t}{\C^n}{\C^n}$。
其标准形由
含有各特征值的 Jordan 块组成。
下图展示了这样一个 $n=6$ 的矩阵，特征值为 $3$、$2$
和~$-1$
（含完整计算的例子在该定理之后）。
\begin{equation*}
  T_3=\begin{mat}
    3  &0  &0  &0 &0 &0 \\
    1  &3  &0  &0 &0 &0  \\
    0  &0  &3  &0 &0 &0  \\
    0  &0  &0  &2 &0 &0  \\
    0  &0  &0  &1 &2 &0  \\
    0  &0  &0  &0 &0 &-1  \\
  \end{mat}
  \qquad
  {\renewcommand{\arraystretch}{1.5}
    \begin{array}{r}
    \vec{\beta}_1\mapsunder{t-3}\vec{\beta}_2\mapsunder{t-3}\zero  \\
    \vec{\beta}_3\mapsunder{t-3}\zero  \\
    \vec{\beta}_4\mapsunder{t-2}\vec{\beta}_5\mapsunder{t-2}\zero  \\
    \vec{\beta}_6\mapsunder{t+1}\zero  \\
    \end{array}}
\end{equation*}
它有四个块，其中两个
与特征值~$3$ 相伴，\( 2 \) 和~\( -1 \) 各有一个。


\begin{theorem}  \label{th:JordanForm}
\index{similar!canonical form}%
\index{canonical form!for similarity}\index{transformation!Jordan form for}
任何变换 $\map{t}{\C^n}{\C^n}$ 都可以表示成
\definend{Jordan 形式},\index{Jordan form!represents similarity classes}\index{Jordan form!definition}
其中每个 $J_{\lambda}$ 都是 Jordan 块。\index{Jordan block}
\begin{equation*}
  \begin{mat}
    J_{\lambda_1}  &            &\text{\textit{--zeroes--}}                 \\
               &J_{\lambda_2}                                              \\
               &     &\ddots                                     \\
             %    &     &                           &J_{\lambda_{k-1}} &     \\
               &     &\text{\textit{--zeroes--}} &            &J_{\lambda_{k}}
  \end{mat}
\end{equation*}
换言之，任何 $\nbyn{n}$ 矩阵都相似于这种形式的某个矩阵。
\end{theorem}

在前面的例子中，若把
$t-3$ 作用于该基，则它会把两个元素
$\vec{\beta}_2$ 和~$\vec{\beta}_3$ 送到~$\zero$。
这提示我们可以用归纳法来做证明，
因为值域空间~$\rangespace{t-3}$ 的维数
小于起始空间的维数。

\begin{proof}
% This proof is closely related to that of
% Five.III.\ref{th:NilMapHasStrBas}, that any nilpotent transformation
% has a basis of disjoint $t$-strings.
%
我们对~$n$ 作归纳。
$n=1$ 的基础情形是平凡的，因为任何 $\nbyn{1}$ 矩阵已经是 Jordan 形式。
归纳步骤假设
$\C^1$, \ldots, $\C^{n-1}$ 上的每个变换都有 Jordan 形式
表示，
并取定 $\map{t}{\C^n}{\C^n}$。

任何变换都至少有一个特征值~$\lambda_1$，因为
它有极小多项式，而
该多项式在复数域上至少有一个根。
该特征值有非零的
特征向量~$\vec{v}$，使得 $(t-\lambda_1)(\vec{v})=\zero$。
% and so the null space of $t-\lambda_1$ is nontrivial.
于是 $\nullspace{t-\lambda_1}$ 的维数，即该映射的零度，
大于零。
由于秩加零度等于定义域的维数，
$t-\lambda_1$ 的秩
严格小于~$n$。
记 $W$ 为值域空间 $\rangespace{t-\lambda_1}$，记
$r$ 为秩，即其维数。

若
$\vec{w}\in W=\rangespace{t-\lambda_1}$，则 $\vec{w}\in\C^n$，于是
$(t-\lambda_1)(\vec{w})\in\rangespace{t-\lambda_1}$。
因此，$t-\lambda_1$ 在~$W$ 上的限制诱导出~$W$ 上的一个变换，
我们仍把它记为 $t-\lambda_1$。
~$W$ 的维数小于~$n$，
所以归纳假设适用，
我们得到~$W$ 的一组由不相交链构成的基，
这些链与 $t-\lambda_1$ 在~$W$ 上
的特征值相伴。
对于~$t-\lambda_1$ 的任一特征值 $\lambda$，
基元素在该变换下的像是
$(t-\lambda_1)-\lambda=t-(\lambda_1-\lambda)$，
我们将把它写作 $t-\lambda_s$。
下面的示意图给出了 $t-\lambda_1$ 的一些链，也就是
$\lambda=0$ 的情形；
若不存在这样的链，论证仍然有效。
\begin{equation*}
\begin{strings}{*{10}{l}}
  \vec{w}_{1,n_1} &\xmapsto{t-\lambda_1} &\cdots
     &\xmapsto{t-\lambda_1} &\vec{w}_{1,1} &\xmapsto{t-\lambda_1} &\zero \\
     \vdotswithin{\xmapsto{t-\lambda_1}}                   \\
  \vec{w}_{q,n_q} &\xmapsto{t-\lambda_1} &\cdots
     &\xmapsto{t-\lambda_1} &\vec{w}_{q,1} &\xmapsto{t-\lambda_1} &\zero \\
  \vec{w}_{q+1,n_{q+1}} &\xmapsto{t-\lambda_2} &\cdots
     &\xmapsto{t-\lambda_2} &\vec{w}_{q+1,1} &\xmapsto{t-\lambda_2} &\zero \\
     \vdotswithin{\xmapsto{t-\lambda_1}}                              \\
  \vec{w}_{k,n_k} &\xmapsto{t-\lambda_i} &\cdots
     &\xmapsto{t-\lambda_i} &\vec{w}_{k,1} &\xmapsto{t-\lambda_i} &\zero \\
\end{strings}
\end{equation*}
这些链的长度可以不同；
例如，也许 $n_1\neq n_k$。

最右边的非 $\zero$ 向量
$\vec{w}_{1,1}$、\ldots、
$\vec{w}_{k,1}$
属于相应
映射的零空间，~$(t-\lambda_1)(\vec{w}_{1,1})=\zero$、\ldots、
$(t-\lambda_i)(\vec{w}_{k,1})=\zero$。
因此，与 $t-\lambda_1$ 相伴的链的个数，
即上图中的 $q$，
等于 $W\cap\nullspace{t-\lambda_1}$ 的维数。

上面的链基是针对空间 $W=\rangespace{t-\lambda_1}$ 的。
我们现在把它扩充为整个~$\C^n$ 的基。
首先，由于
向量 $\vec{w}_{1,n_1}$、\ldots、~$\vec{w}_{q,n_q}$
中的每一个都是~$\rangespace{t-\lambda_1}$ 的元素，每一个都是来自~$\C^n$ 的
某个向量的像。
给每条链前面添加这样一个向量，
\( \vec{x}_1 \)、\ldots、\( \vec{x}_q \)，
如下所示。
其次，$W\cap\nullspace{t-\lambda_1}$ 的维数是~$q$，
而~$W$ 的维数是~$r$，
所以按账目计算，需要有一个维数为 $r-q$ 的子空间
$Y\subseteq\nullspace{t-\lambda_1}$，
它与 $W$ 的交是~$\set{\zero}$。
于是，
取 $r-q$ 个线性无关的
向量 $\vec{y}_1$、\ldots、$\vec{y}_{r-q}\in\nullspace{t-\lambda_1}$
并把它们并入各链的名单之中，同样如下所示。
\begin{equation*}
\begin{strings}{*{10}{c}}
  \vec{x}_1  &\xmapsto{t-\lambda_1}
     &\vec{w}_{1,n_1} &\xmapsto{t-\lambda_1} &\cdots
     &\xmapsto{t-\lambda_1} &\vec{w}_{1,1} &\xmapsto{t-\lambda_1} &\zero \\
     \vdotswithin{\xmapsto{t-\lambda_1}}                   \\
  \vec{x}_q  &\xmapsto{t-\lambda_1}
     &\vec{w}_{q,n_q} &\xmapsto{t-\lambda_1} &\cdots
     &\xmapsto{t-\lambda_1} &\vec{w}_{q,1} &\xmapsto{t-\lambda_1} &\zero \\
  &&\vec{w}_{q+1,n_{q+1}} &\xmapsto{t-\lambda_2} &\cdots
     &\xmapsto{t-\lambda_2} &\vec{w}_{q+1,1} &\xmapsto{t-\lambda_2} &\zero \\
     &&\vdotswithin{\xmapsto{t-\lambda_1}}                              \\
  &&\vec{w}_{k,n_k} &\xmapsto{t-\lambda_i} &\cdots
     &\xmapsto{t-\lambda_i} &\vec{w}_{k,1} &\xmapsto{t-\lambda_i} &\zero \\
     &&&&&&\vec{y}_1  &\xmapsto{t-\lambda_1} &\zero \\
     &&&&&&\vdotswithin{\xmapsto{t-\lambda_1}}                              \\
     &&&&&&\vec{y}_{r-q}  &\xmapsto{t-\lambda_1} &\zero \\
\end{strings}
\end{equation*}
我们将证明这就是所想要的~$\C^n$ 的基。

由于按账目计算，该集合中向量的个数等于空间的维数，
所以只要验证其元素线性无关，我们就完成了。
假设下式等于~$\zero$。
\begin{equation*}
  a_1\vec{x}_1+\cdots+a_q\vec{x}_q
    +b_{1,n_1}\vec{w}_{1,n_1}+\cdots+b_{k,1}\vec{w}_{k,1}
    +c_1\vec{y}_1+\cdots+c_{r-q}\vec{y}_{r-q}
  \tag{$*$}
\end{equation*}
我们首先证明，由此可推出所有的 $a$ 都为零。
把 $t-\lambda_1$ 作用于~($*$)。
向量 $\vec{y}_1$、\ldots、$\vec{y}_{r-q}$
被送到~$\zero$。
每个向量 $\vec{x}_i$ 都被送到
$\vec{w}_{i,n_i}$。
至于各个 $\vec{w}$，它们具有如下性质：对某个~$\lambda_k$，
$(t-\lambda_k)\,\vec{w}_{i,j}=\vec{w}_{i,j-1}$。
记 $\hat{\lambda}$ 为 $\lambda_1-\lambda_k$，便得到
\begin{equation*}
  (t-\lambda_1)(\vec{w}_{i,j})
    =(t-(\lambda_1-\hat{\lambda})+\hat{\lambda})(\vec{w}_{i,j})
    =\vec{w}_{i,j-1}+\hat{\lambda}\cdot\vec{w}_{i,j}
  \tag{$**$}
\end{equation*}
于是，把 $t-\lambda_1$ 作用于~($*$) 之后最终得到的是 $\vec{w}$ 的一个
线性组合。
这些向量线性无关，因为
它们构成~$W$ 的基，所以在这个线性组合中
系数是~$0$。
要完成 $a_i=0$（$i=1,\ldots,q$）为零的证明，
只需验证：在把 $t-\lambda_1$ 作用于~($*$) 之后，
$\vec{w}_{i,n_i}$ 的系数是 $a_i$。
而这确实成立，因为对 $i=1,\ldots,q$，
方程~($**$) 中 $\hat{\lambda}=0$，
所以 $\vec{w}_{i,n_i}$ 仅是 $x_{i}$ 的像。
% Thus $a_1$, \ldots, $a_q$ are zero.

各 $a$ 等于零之后，~($*$) 剩下的部分就是
\begin{equation*}
  \zero
  =
    b_{1,n_1}\vec{w}_{1,n_1}+\cdots+b_{k,1}\vec{w}_{k,1}
    +c_1\vec{y}_1+\cdots+c_{r-q}\vec{y}_{r-q}
\end{equation*}
改写为 $-(b_{1,n_1}\vec{w}_{1,n_1}+\cdots+b_{k,1}\vec{w}_{k,1})=
    c_1\vec{y}_1+\cdots+c_{r-q}\vec{y}_{r-q}$。
各 $\vec{w}$ 来自~$W$，而各 $\vec{y}$ 来自~$Y$，且
这两个空间只有平凡的交。
因此 $b_{1,n_1}\vec{w}_{1,n_1}+\cdots+b_{k,1}\vec{w}_{k,1}=\zero$，
且 $c_1\vec{y}_1+\cdots+c_{r-q}\vec{y}_{r-q}=\zero$，
从而各 $b$ 与各 $c$ 全都为零。
\end{proof}





% Recall that Lemma~III.\ref{le:RangeAndNullChains} shows that for any
% transformation~$t$ on an $n$~dimensional space
% the range spaces of iterates are stable
% \begin{equation*}
%   \rangespace{t^n}=\rangespace{t^{n+1}}=\cdots=\genrangespace{t}
% \end{equation*}
% as are the null spaces.
% \begin{equation*}
%   \nullspace{t^n}=\nullspace{t^{n+1}}=\cdots=\gennullspace{t}
% \end{equation*}
% Thus,
% the generalized null space $\gennullspace{t}$ and the generalized range space
% $\genrangespace{t}$ are $t$~invariant.
% % For the generalized null space, if $\vec{v}\in\gennullspace{t}$ then
% % $t^n(\vec{v})=\zero$ where $n$ is the dimension of the underlying space
% % and so $t(\vec{v})\in\gennullspace{t}$ because
% % $t^n(\,t(\vec{v})\,)$ is zero also.
% % For the generalized rangespace, if $\vec{v}\in\genrangespace{t}$ then
% % $\vec{v}=t^n(\vec{w})$ for some $\vec{w}$ and then
% % $t(\vec{v})=t^{n+1}(\vec{w})=t^n(\,t(\vec{w})\,)$
% % shows that $t(\vec{v})$ is also a member of $\genrangespace{t}$.
% Also,
% $\gennullspace{t-\lambda_i}$ and $\genrangespace{t-\lambda_i}$
% are $t-\lambda_i$~invariant.

